\section{The shared contraction interface and what is imported}
\label{sec:interface}

Write
\[
 \lambda=\lceil\log_2 s\rceil,\qquad E(s)=s/\sqrt\lambda,
 \qquad \kappa=\pi\sqrt2/\ln2.
\]
The comparison problem is a Boolean tensor network with at most $s$ wires,
maximum initial degree at most $\lceil s^{9/10}\rceil$, and initial
precision and coefficient truncation length at most $s$. Its entries lie
in the power-of-two coefficient rings of the OpenAI construction. Original
entry routines have the prescribed polynomial-description and
$2^{O(s^{9/10})}$-time interface. Retained tensors are evaluated at their
immutable creation precision and then lifted as specified by the source.
These premises matter: an arbitrary exponentially expensive node oracle
would not satisfy the theorem.

We import the source's exact deformation, recovery identity, child-network
construction, SAT reductions, and their correctness
\cite{openai,contraction}. We also import nearly lossless asymmetric packet
packing, the sharp tapered precision, and full-block Cauchy modes from the
earlier refinements \cite{asymmetric,precision,joint,adaptive,tree}.
This report proves a new method for evaluating the resulting finite mode
graph. It does not replace the proofs of the imported identities.

\subsection{The exact mode graph}

At an active call, let $M$ be its number of wires, $u$ its retained
valuation range, and $r$ its fixed alphabet size. After choosing an exact
pair-weight allocation, the local answer has the form
\begin{equation}\label{eq:modegraph}
 \sum_{\xi_e\in\mathcal I_e}\prod_{g=1}^{r}
    Z_g\bigl((\xi_e)_{e\ni g}\bigr).
\end{equation}
Pairs $e=(g,h)$ are directed and have distinct endpoints. The list
$\mathcal I_e$ contains the feasible mode-index tuples of the allocated
exact weight; $Z_g$ is an actual child-network contraction. Each child has
at most $(2/r+o(1))M$ Boolean wires. Empty lists make the allocation zero.

Because $r$ is fixed and all budgets are polynomial in $s$, the number of
exact pair-weight allocations is polynomial. The usual integer knapsack
recurrence computes each list size and supplies ranking and unranking with
polynomial overhead. This is not a unit-cost counting oracle.

Put $h_e=\log_2|\mathcal I_e|$ and
$\K(M,u)=\kappa M/\sqrt u$. The imported common-tilt estimate is
\begin{equation}\label{eq:entropy}
 \sum_e z_e h_e\le
  \left(\frac{\|z\|_2}{r}+o(1)\right)\K(M,u)
 \qquad (z_e\ge0\hbox{ fixed}).
\end{equation}
The error is uniform over the allocations and all active calls of the
admissible schedules below, for each fixed finite family of $z$.

For clarity, its leading constants can be recovered directly. The
full-Cauchy weights may be bounded below by
$w_b=1+4b-2\operatorname{popcount}(b)$. Their logarithmic generating
function is
\[
 \sum_{b\ge0}-\ln(1-e^{-\theta w_b})
 =\frac{\pi^2}{24\theta}+O\bigl(\log^2(2/\theta)\bigr).
\]
The color multiplicity on one directed pair is
$(2+o(1))M/(r^2D)$, and the working precision is
$(6+o(1))MD/u$. Applying the generating-function bound with tilt
$\theta/z_e$, summing with weight $z_e$, and optimizing the leading terms
gives $\pi\sqrt2 M\|z\|_2/(r\sqrt u)$ in natural logarithms.
Division by $\ln2$ gives \eqref{eq:entropy}.

\subsection{Paid recursive calls}

The first stage caches each $Z_g$ on all incident pair indices. Its
incident set has $2(r-1)$ directed pairs, so the total number of child
calls has exponent at most $(\alpha_r+o(1))\K$, where
\begin{equation}\label{eq:alpha}
 \alpha_r=\frac{\sqrt{2(r-1)}}r.
\end{equation}
A local contraction plan then combines these cached functions. If its
certified nonrecursive time coefficient is $\beta_r$, with backend slack
$\delta>0$ charged explicitly, the recurrence is
\begin{equation}\label{eq:paid-recurrence}
 T_d\le\poly(s)\left[
     2^{(\alpha_{r_d}+o(1))\K_d}T_{d+1}
       +2^{((1+\delta)\beta_{r_d}+o(1))\K_d}\right].
\end{equation}
The local space coefficient is at most $\beta_{r_d}$; the input/output
term of the matrix cost pays for every table. Recursive tables on a stack
add in space, whereas ancestor child-call exponents add in time. Ignoring
that distinction would give an incorrect headline bound.

\subsection{What an improved wire coefficient means for SAT}

All numerical coefficients in this report multiply $E(s)$ per
\emph{normalized Boolean wire}. The known reduction yields, for fixed $k$,
the inherited density-sensitive order
\begin{equation}\label{eq:sat-envelope}
 \poly_k(n,m)\,
 2^{O_k\left(\frac{n}{\sqrt{\log n}}
       \min\{\log(2+m/n),\log\log n\}\right)}.
\end{equation}
The direct clause-count branch and all stronger-assumption low-degree
branches remain available. Smaller explicit wire coefficients give smaller
certified contraction upper bounds. They do not determine numerical
coefficients per SAT variable without the reduction's wire constant, and
they do not change the logarithmic order in \eqref{eq:sat-envelope}.
